%!TeX program = xelatex
\PassOptionsToPackage{quiet}{fontspec}
\PassOptionsToPackage{fontset=fandol}{ctex}
\PassOptionsToPackage{table}{xcolor}
\documentclass[10pt,compsoc,UTF8]{CjC}
\usepackage{thesis}
\usepackage{tabularx,longtable,hyperref,attachfile2}
\usepackage{etoolbox}
\makeatletter
\patchcmd{\l@section}{1.5em}{2.5em}{}{\PackageError{report}{TOC number width patch failed}{Check CjC template.}}
\makeatother
\setcitestyle{numbers,square}
\hypersetup{colorlinks=true,linkcolor=black,urlcolor=blue!55!black,
pdftitle={基于量价因子的十组回测与XGBoost选股},pdfauthor={岳瀚阳}}
\renewcommand{\today}{2026年11月16日}
\mycourse{《量化投资》}
\mytitle{基于量价因子的十组回测与XGBoost选股}
\myauthor{岳瀚阳}\mystudentid{3230104443}\myschool{管理学院}\myteacher{曾涛、彭旺}
\setlength{\emergencystretch}{2em}
\setcounter{tocdepth}{1}
\newcommand{\resultfigure}[2]{\begin{figure}[H]\centering
\includegraphics[width=.96\linewidth,height=.70\textheight,keepaspectratio]{#1.pdf}
\caption{#2。上图为十个独立多头组，G10为看好组，红色虚线为同池全体等权费后基准；中图为每腿100\%目标的费后理论多空；下图为G10净值除以基准净值。}\end{figure}}
\newcommand{\Searches}{24}
\newcommand{\CVFits}{3456}
\newcommand{\CoarseFits}{2304}
\newcommand{\FineFits}{1152}
\newcommand{\Models}{6}
\newcommand{\ValidationModels}{18}
\newcommand{\Series}{672}
\newcommand{\Trades}{2,406,703}
\newcommand{\Tests}{93}
\newcommand{\BestModelPool}{中证500}
\newcommand{\BestModelPeriod}{10日}
\newcommand{\BestModelLS}{22.19\%}
\newcommand{\BestModelLSSharpe}{2.55}
\newcommand{\BestModelLSDrawdown}{-3.68\%}
\newcommand{\PlotPool}{中证500}
\newcommand{\PlotPeriod}{10日}
\newcommand{\PlotMeanTotal}{9.14\%}

\begin{document}
\raggedbottom
\setpage{}
\setcounter{page}{2}
\pagestyle{fancy}\fancyhf{}
\fancyhead[L]{\zihao{-5}\themyauthor\quad\themystudentid}
\fancyhead[R]{\zihao{-5}\themycourse}
\fancyfoot[C]{\zihao{-5}\thepage}
\setlength{\headheight}{18pt}
\Title{\themytitle}{\themyauthor}{\themystudentid}
\myabstract{
本文按固定时间划分开展量价因子选股实验：2016--2017年用于训练，2018年用于回测。
在历史沪深300、中证500及本地覆盖的全部A股三个股票池中构造9个候选因子，
执行十组多头与理论多空回测，并以同池全体合格股票的等权组合作为市场参照。
信号在T日收盘后形成，T+1开盘执行；买卖各收0.0015，完整往返0.003。

XGBoost分别预测5、10、20交易日连续开盘收益。因子去重、择优、调参及选池、选期限均只使用训练期数据。
每池每期限搜索96组候选，前16组晋级2017年前三个季度的三折时间验证。
最后使用2016--2017年已实现标签拟合\Models 个固定模型，2018全年不重训、不重新选参。
总计\Searches 次搜索、\CVFits 次验证拟合，另有\ValidationModels 个训练期验证重拟合。
按训练期验证确定的XGBoost组合为\BestModelPool、\BestModelPeriod，
2018年费后多头累计收益\BestModelLong，同池等权\BestModelEW，相对净值收益\BestModelRelative。
报告完整列出十分组收益、三池结果、市场月度变化与费用反事实；2018年收益不反向用于选择展示组合。
}
\keyword{量价因子；固定训练；XGBoost；2018回测；十组收益；等权基准}
\tableofcontents

\clearpage
\section{课题具体说明}
研究基于原作业和本地研报，在2016--2017年完成选型与训练，于2018年统一测试。
九候选包括聪明钱、CPV、CDPDP$^{\prime}$、理想振幅、理想反转、CDPP、四象限CDPDP、低振幅长端动量和分钟偏度。
前八项来自研报定义或其原始引用，分钟偏度为原作业扩展；不纳入尾部收益因子。

单因子和模型均采用十等分组，G10统一为看好组，G1为看空组。
因子方向预先确定：低振幅长端动量高值做多，其余因子低值做多，模型为高预测值做多。
每个多头组初始净资产为1、目标多头100\%；多空组合多、空各100\%，总名义敞口200\%，
独立维护持仓、现金及实际费用，不能把十组收益简单相减视为已经完成成交的策略。

沪深300、中证500使用当时历史成分；全部A股限于本地日行情覆盖的A股证券。
每项策略有三池$\times$三个固定期限共9个候选，选图组合由训练期时间验证决定。
固定20交易日不称作自然月。2018年仅做一次冻结后的评估，不用其收益挑股票池、期限、因子方向或参数。

\section{数据描述与交易处理}
分钟文件覆盖2016--2018年731个交易日，日行情额外提供200个交易日预热。
个股分钟和滚动窗口可复用，截面去极值、中性化和标准化在各股票池内分别计算。
因子及动态成分在2018年仍按当日信息更新；冻结的是模型权重、入选因子、参数和选池/期限规则。
模型使用原生缺失值处理，评分入口只要求当日有至少3个候选因子，不查看未来标签是否完整。

\begin{table}[H]\centering\normalsize
\caption{数据用途与时间对齐}
\begin{tabularx}{\linewidth}{p{28mm}X}\toprule
数据 & 处理方式\\\midrule
分钟行情 & 分钟开高低收、成交量、真实成交额及成交笔数；构造10日、20日信号。\\
官方复权 & 使用当日可知因子形成连续价格；收益及排序统一单位，实际限价判断仍用原价。\\
历史ST & T日剔除ST信号；执行日ST仅限制风险增加，减仓与回补仍允许。\\
停复牌 & 按09:30时段判断开盘成交，不使用稍后复牌价或盘中高低点倒推早盘成交。\\
实际限价 & 开盘等于上限不买、等于下限不卖；缺失或与报价冲突时阻止可疑成交。\\
成分、市值与资格 & T日历史成分、市值和20日风格变量；上市满60交易日以全历史首次正报价近似，剔除当日ST与停牌信号。\\\bottomrule
\end{tabularx}\end{table}
受阻买入留现金，受阻退出保留并逐日重试；未知行情沿用最近估值并标记。
2017年最后一个交易日收盘生成首次测试信号，2018年首个交易日开盘建仓，末日开盘尝试退出。
未另计借券、融资和冲击成本；理论多空依赖可借券及融资假设。

\clearpage
\section{因子构造逻辑}
\begin{table}[H]\centering\normalsize
\caption{九个候选因子的定义}
\begin{tabularx}{\linewidth}{p{37mm}X}\toprule
因子 & 构造方法\\\midrule
SmartMoney & 对过去10个交易日分钟数据整体按$|R|/V^{0.25}$排序，取累计20\%成交量，边界按比例截取。
以真实成交额计算聪明钱VWAP与全部VWAP之比；跨除权日统一价格/份额单位。\\
CPV & 对日内价格与成交量相关系数，取20日均值、标准差和时间趋势。
均值/标准差先做市值处理，再各自剔除反转并标准化；趋势剔除市值、反转、波动和换手，最后分阶段合成。\\
CDPDP$^{\prime}$ & 对相邻分钟价格差分按同涨、同跌分别求相关，20日聚合、市值处理、子项分别Z后相加。\\
IdealAmplitude & 20日窗口剔除停牌和一字涨跌停，有效日至少10天。
按连续收盘价取最高、最低各25\%有效日，分别计算$H/L-1$均值后相减。每端日数向下取整。\\
IdealReversal & 日总成交额/总成交笔数；20日中按平均单笔金额取高10日和低10日，各自加总对应日收益后相减。\\
CDPP & 计算$\Delta P_t$与$P_{t+1}$在正、负差分条件下的相关，20日聚合、市值处理、子项各自Z后相减。\\
CDPDP-4 & 四象限相关子项分别处理，采用$Z(++)-Z(+-)-Z(-+)+Z(--)$。\\
LowAmpMomentum & 过去160个交易日中取振幅最低70\%的112日，将对应日收益加总，高值做多。\\
Skewness & 计算日内相邻分钟收益的三阶标准化中心矩，再取20日均值；属于原作业扩展。\\\bottomrule
\end{tabularx}\end{table}

各因子按其要求截面去极值、标准化和控制风格，未统一声称已经剔除完整Barra暴露。
相关子序列不足3对或零方差置0，整日缺行情保持缺失；分钟首笔、边界插值和并列价格处理属于公开的实现约定。
完整Barra风格、逐笔大小单方向及可借券供应未在本地数据中提供，因此不将本研究称为券商私有代码和数据的逐字复制。



\clearpage
\section{固定训练、时间验证与2018冻结回测}
\begin{table}[H]\centering\normalsize
\caption{严格分开的训练、验证与测试范围}
\begin{tabularx}{\linewidth}{p{35mm}X}\toprule
环节 & 时间和信息约束\\\midrule
最终训练 & 2016--2017年信号，且标签退出日必须早于2018年1月1日；生成训练标签的行情本身截止2017年。\\
内层第1折 & 训练标签在2017年1月1日前实现，验证2017年第一季度。\\
内层第2折 & 训练标签在2017年4月1日前实现，验证2017年第二季度。\\
内层第3折 & 训练标签在2017年7月1日前实现，验证2017年第三季度。\\
选池及期限 & 使用上述时间验证预测回测的G10费后年化；不读取2018收益。\\
2018测试 & 最终9个模型一次拟合后全年冻结，无重训、无测试期早停、无测试期重新去重择优。\\\bottomrule
\end{tabularx}\end{table}
标签为$O^{adj}_{t+N+1}/O^{adj}_{t+1}-1$的截面去均值，$N$分别为5、10、20。
跨入2018年才实现的标签不进入训练；例如年末信号即使发生在2017年，也不能因其信号日较早而进入训练。
2017年第四季度仍进入最终拟合；参数验证窗口固定为前三季度，便于三个期限使用同一验证回测区间。

验证阶段每折只用该季度开始前已实现的训练标签拟合；季度衔接时，用前一交易日收盘信号在下一交易日开盘执行。
验证预测与真实交易回测均保留当时可计算但未来标签缺失的股票，不以未来缺失过滤预测入口。
参数选择及选池、选期限共用训练期验证，因此验证收益只用于选型，不能作为独立测试证据；2018仅用于本轮冻结后的测试评价。

\clearpage
\section{相关性去重与训练期因子择优}
相关性使用训练期逐日截面Spearman相关并等权平均；质量为日RankIC绝对均值除以标准差。
至少20个有效IC日、训练行覆盖不少于50\%，按质量顺序择优；与已选代表的相关系数绝对值达到阈值即剔除。
阈值0.70/0.80/0.90和最多4/6/9项的预算与树模型参数一起搜索。
每个内层折在自己的训练部分重算筛选；最终模型在2016--2017年完整已实现样本上重新筛选并拟合一次。

\begin{figure}[H]\centering
\includegraphics[width=.82\linewidth,height=.40\textheight,keepaspectratio]{correlation.pdf}
\caption{训练期所选模型（\BestModelPool、\BestModelPeriod）的2016--2017年相关矩阵}
\end{figure}
\begin{table}[H]\centering\normalsize
\caption{最终固定模型中的因子入选次数：每个股票池3个期限}
\begin{tabular}{lrrr}\toprule
候选因子 & 沪深300 & 中证500 & 全部A股\\\midrule
SmartMoney & 3 & 1 & 3 \\
CPV & 3 & 3 & 3 \\
CDPDP$^{\prime}$ & 2 & 1 & 3 \\
IdealAmplitude & 2 & 1 & 1 \\
IdealReversal & 3 & 3 & 3 \\
CDPP & 3 & 3 & 3 \\
CDPDP-4 & 2 & 1 & 0 \\
LowAmpMomentum & 3 & 1 & 1 \\
Skewness & 3 & 3 & 3 \\
\bottomrule
\end{tabular}\end{table}

\clearpage
\section{XGBoost充分调参与最终参数}
每池每期限进行一次完整搜索：96组初筛，前16组晋级三折时间验证。
初筛最多300树、25轮早停；精筛最多1200树、50轮早停。
早停仅监控训练期验证的加权RMSE，按三折日RankIC均值减0.1倍折间标准差择优。
最终树数为获胜候选在三折的最优树数中位数；2018测试过程不设置评估集，也不调整树数。

\begin{table}[H]\centering\normalsize
\caption{固定实验的搜索空间}
\begin{tabularx}{\linewidth}{p{44mm}X}\toprule
参数 & 搜索范围\\\midrule
max\_depth & 2、3、4、5、6\\
learning\_rate & 对数均匀0.015--0.20\\
min\_child\_weight & 1、3、10、30、60、100\\
subsample / colsample & 各自均匀0.55--1.00\\
reg\_alpha / reg\_lambda & $\alpha=0$或$10^{-4}$至约2；$\lambda=0.1$至100\\
gamma / max\_bin & $\gamma=0$或$10^{-5}$至0.1；128或256箱\\
相关阈值 / 因子预算 & 0.70、0.80、0.90；最多4、6、9项\\\bottomrule
\end{tabularx}\end{table}
日期总训练权重相同，避免股票较多的日期主导损失；随机种子固定。
共完成\Searches 次搜索、\CVFits 次验证拟合；为计算训练期策略验证收益另拟合\ValidationModels 个折内模型，
最后拟合\Models 个固定模型用于2018全年。候选与验证边界均保留在各池models子目录。

\begin{table}[H]\centering\normalsize
\caption{2016--2017训练期选出的九个最终配置}
\setlength{\tabcolsep}{4pt}
\begin{tabular}{llrrrrrr}\toprule
股票池 & 期限 & 深度 & 学习率 & 树数 & 相关阈值 & 因子上限 & 验证RankIC\\\midrule
沪深300 & 5日 & 3 & 0.050 & 47 & 0.90 & 6 & 0.082 \\
沪深300 & 10日 & 2 & 0.198 & 13 & 0.90 & 9 & 0.104 \\
沪深300 & 20日 & 2 & 0.184 & 19 & 0.90 & 9 & 0.133 \\
中证500 & 5日 & 3 & 0.026 & 54 & 0.80 & 4 & 0.089 \\
中证500 & 10日 & 2 & 0.132 & 12 & 0.90 & 9 & 0.106 \\
中证500 & 20日 & 4 & 0.034 & 55 & 0.70 & 4 & 0.105 \\
全部A股 & 5日 & 5 & 0.018 & 199 & 0.70 & 9 & 0.114 \\
全部A股 & 10日 & 5 & 0.075 & 40 & 0.70 & 6 & 0.132 \\
全部A股 & 20日 & 6 & 0.017 & 146 & 0.70 & 6 & 0.144 \\
\bottomrule
\end{tabular}\end{table}

\clearpage
\section{训练期选池与2018三池比较}
对九个因子及XGBoost，分别按2017年前三季度时间验证预测的G10费后年化，在三池三期限中择优。
该规则在查看2018结果之前固定。下面每池一列也先在训练期验证中选出该池的期限，再列对应的2018累计收益；
不是把每池2018年表现最高的期限拿来汇报。

\begin{table}[H]\centering\normalsize
\caption{三池比较：训练期验证选期限后的2018费后累计收益}
\setlength{\tabcolsep}{4pt}
\begin{tabular}{lrrrl}\toprule
策略 & 沪深300 & 中证500 & 全部A股 & 训练期选图池/期限\\\midrule
SmartMoney & -34.79\% & -34.62\% & -29.84\% & 沪深300/5日 \\
CPV & -22.94\% & -29.11\% & -28.15\% & 沪深300/20日 \\
CDPDP$^{\prime}$ & -24.95\% & -25.98\% & -26.20\% & 沪深300/5日 \\
IdealAmplitude & -24.34\% & -33.16\% & -27.86\% & 沪深300/20日 \\
IdealReversal & -23.21\% & -31.29\% & -25.20\% & 中证500/10日 \\
CDPP & -19.44\% & -32.13\% & -20.78\% & 沪深300/20日 \\
CDPDP-4 & -25.47\% & -26.16\% & -26.38\% & 沪深300/20日 \\
LowAmpMomentum & -25.35\% & -26.92\% & -30.56\% & 沪深300/5日 \\
Skewness & -23.68\% & -30.09\% & -27.53\% & 沪深300/20日 \\
XGBoost & -17.03\% & -32.13\% & -21.67\% & 沪深300/20日 \\
\bottomrule
\end{tabular}\end{table}

\begin{table}[H]\centering\normalsize
\caption{训练期选定的十项组合：验证与2018测试结果}
\setlength{\tabcolsep}{3pt}
\begin{tabular}{llrrrrr}\toprule
策略 & 股票池/期限 & 验证年化 & 2018多头 & 池等权 & 相对收益 & 多空收益\\\midrule
SmartMoney & 沪深300/5日 & 6.66\% & -34.79\% & -27.15\% & -10.48\% & -11.56\% \\
CPV & 沪深300/20日 & 27.68\% & -22.94\% & -27.01\% & 5.59\% & 10.57\% \\
CDPDP$^{\prime}$ & 沪深300/5日 & 14.66\% & -24.95\% & -27.15\% & 3.02\% & 4.17\% \\
IdealAmplitude & 沪深300/20日 & 15.17\% & -24.34\% & -27.01\% & 3.67\% & 12.35\% \\
IdealReversal & 中证500/10日 & 7.34\% & -31.29\% & -32.09\% & 1.18\% & 2.11\% \\
CDPP & 沪深300/20日 & 21.38\% & -19.44\% & -27.01\% & 10.37\% & 31.24\% \\
CDPDP-4 & 沪深300/20日 & 18.07\% & -25.47\% & -27.01\% & 2.12\% & 2.93\% \\
LowAmpMomentum & 沪深300/5日 & 35.76\% & -25.35\% & -27.15\% & 2.46\% & 3.21\% \\
Skewness & 沪深300/20日 & 12.13\% & -23.68\% & -27.01\% & 4.56\% & 13.40\% \\
XGBoost & 沪深300/20日 & 24.19\% & -17.03\% & -27.01\% & 13.68\% & 29.79\% \\
\bottomrule
\end{tabular}\end{table}
除验证列为年化外，其余收益列均为2018累计收益。
相对收益为$\mathrm{NAV}_{G10}/\mathrm{NAV}_{EW}-1$，不是两条累计收益直接相减。
2018年可能存在表现更高的其他候选，仍不以其结果反向更换图中组合。

\clearpage
\section{九个固定XGBoost模型的2018回测}
下面完整列出三个股票池和三个期限，便于比较；每个模型2018年都使用2017年底冻结的权重和入选因子。
G10多头及各组均有独立初始资金1；理论多空为多、空各100\%，资本敞口不同。

\begin{table}[H]\centering\normalsize
\caption{固定模型：训练期验证与2018多头、多空收益}
\setlength{\tabcolsep}{3pt}
\begin{tabular}{llrrrrrr}\toprule
股票池 & 期限 & 验证年化 & 多头累计 & 多头年化 & 池等权累计 & 多空累计 & 多头夏普\\\midrule
沪深300 & 10日 & 19.16\% & -23.13\% & -23.71\% & -27.17\% & 20.44\% & -1.05 \\
沪深300 & 20日 & 24.19\% & -17.03\% & -17.47\% & -27.01\% & 29.79\% & -0.71 \\
沪深300 & 5日 & 19.02\% & -27.05\% & -27.71\% & -27.15\% & 17.32\% & -1.30 \\
中证500 & 10日 & 3.90\% & -32.32\% & -33.08\% & -32.09\% & 17.95\% & -1.55 \\
中证500 & 20日 & 3.99\% & -29.40\% & -30.10\% & -32.10\% & 19.70\% & -1.44 \\
中证500 & 5日 & 5.95\% & -32.13\% & -32.89\% & -32.10\% & 19.02\% & -1.55 \\
全部A股 & 10日 & -5.21\% & -21.28\% & -21.82\% & -30.37\% & 51.71\% & -0.89 \\
全部A股 & 20日 & -7.10\% & -22.16\% & -22.72\% & -30.52\% & 35.74\% & -0.96 \\
全部A股 & 5日 & -3.80\% & -21.67\% & -22.21\% & -30.15\% & 59.58\% & -0.91 \\
\bottomrule
\end{tabular}\end{table}

\begin{table}[H]\centering\normalsize
\caption{预测覆盖与全体等权基准：每池按训练期选定的模型期限}
\begin{tabular}{llrrr}\toprule
股票池 & 期限 & 平均评分股票数 & 平均基准股票数 & 数量比\\\midrule
沪深300 & 20日 & 284 & 289 & 98.26\% \\
中证500 & 5日 & 461 & 474 & 97.32\% \\
全部A股 & 5日 & 3162 & 3258 & 97.04\% \\
\bottomrule
\end{tabular}\end{table}
基准在当日股票池全体合格股票中等权，不按因子缺失或模型预测过滤；策略受分钟数据及因子可计算性限制。
它是同池的等权市场参照，不是官方市值加权指数，覆盖差异也会影响相对收益。
两者使用相同调仓、开盘限制、ST、停复牌、现金处理和每边0.0015费用。

\clearpage
\subsection{所选组合的年化收益与夏普}
下表统一按2018实际日收益计算，列出训练期所选九因子及模型的收益和风险。多头与理论多空敞口不同，需结合各自口径阅读。

\begin{table}[H]\centering\normalsize
\caption{2018年所选组合的年化收益与夏普}
\setlength{\tabcolsep}{4pt}
\begin{tabular}{llrrrr}\toprule
策略 & 股票池/期限 & 多头年化 & 多头夏普 & 多空年化 & 多空夏普\\\midrule
SmartMoney & 沪深300/5日 & -35.58\% & -1.66 & -11.88\% & -0.81 \\
CPV & 沪深300/20日 & -23.51\% & -1.19 & 10.89\% & 0.96 \\
CDPDP$^{\prime}$ & 沪深300/5日 & -25.56\% & -1.59 & 4.29\% & 0.36 \\
IdealAmplitude & 沪深300/20日 & -24.94\% & -0.98 & 12.73\% & 0.90 \\
IdealReversal & 中证500/10日 & -32.03\% & -1.33 & 2.17\% & 0.25 \\
CDPP & 沪深300/20日 & -19.94\% & -0.92 & 32.27\% & 2.48 \\
CDPDP-4 & 沪深300/20日 & -26.10\% & -1.55 & 3.01\% & 0.29 \\
LowAmpMomentum & 沪深300/5日 & -25.98\% & -1.03 & 3.31\% & 0.27 \\
Skewness & 沪深300/20日 & -24.28\% & -1.23 & 13.82\% & 1.15 \\
XGBoost & 沪深300/20日 & -17.47\% & -0.71 & 30.77\% & 1.88 \\
\bottomrule
\end{tabular}\end{table}

\clearpage
\section{十个多头组分别赚亏多少}
以下数字统一为\textbf{2018年费后累计收益}，G1至G10都是独立的100\%目标多头组合。
十列拆成两表，保持正文五号字。先列中证500固定5日参照，再列训练期验证确定的各策略组合。

\begin{table}[H]\centering\normalsize
\caption{中证5005日参照：2018年G1至G5多头收益}
\setlength{\tabcolsep}{3pt}
\begin{tabular}{llrrrrr}\toprule
策略 & 股票池/期限 & G1 & G2 & G3 & G4 & G5\\\midrule
SmartMoney & 中证500/5日 & -52.73\% & -42.48\% & -42.46\% & -33.99\% & -37.50\% \\
CPV & 中证500/5日 & -51.00\% & -45.84\% & -43.28\% & -39.17\% & -40.13\% \\
CDPDP$^{\prime}$ & 中证500/5日 & -46.53\% & -36.62\% & -36.88\% & -41.22\% & -41.34\% \\
IdealAmplitude & 中证500/5日 & -52.18\% & -45.99\% & -38.60\% & -30.92\% & -34.55\% \\
IdealReversal & 中证500/5日 & -40.74\% & -41.90\% & -40.58\% & -43.98\% & -35.47\% \\
CDPP & 中证500/5日 & -46.12\% & -46.64\% & -43.20\% & -42.23\% & -36.55\% \\
CDPDP-4 & 中证500/5日 & -42.08\% & -40.24\% & -37.64\% & -39.78\% & -34.77\% \\
LowAmpMomentum & 中证500/5日 & -37.67\% & -37.15\% & -38.11\% & -36.40\% & -32.80\% \\
Skewness & 中证500/5日 & -42.34\% & -42.35\% & -38.67\% & -39.39\% & -39.05\% \\
XGBoost & 中证500/5日 & -52.21\% & -49.41\% & -38.30\% & -38.55\% & -36.61\% \\
\bottomrule
\end{tabular}\end{table}
\begin{table}[H]\centering\normalsize
\caption{中证5005日参照：2018年G6至G10与全体等权收益}
\setlength{\tabcolsep}{3pt}
\begin{tabular}{llrrrrrr}\toprule
策略 & 股票池/期限 & G6 & G7 & G8 & G9 & G10 & 池等权\\\midrule
SmartMoney & 中证500/5日 & -37.50\% & -35.30\% & -38.24\% & -35.70\% & -33.60\% & -32.10\% \\
CPV & 中证500/5日 & -35.61\% & -36.01\% & -33.73\% & -33.82\% & -27.63\% & -32.10\% \\
CDPDP$^{\prime}$ & 中证500/5日 & -38.79\% & -35.60\% & -34.17\% & -34.82\% & -26.84\% & -32.10\% \\
IdealAmplitude & 中证500/5日 & -40.92\% & -36.16\% & -38.89\% & -36.63\% & -27.67\% & -32.10\% \\
IdealReversal & 中证500/5日 & -38.23\% & -38.47\% & -36.92\% & -36.78\% & -33.25\% & -32.10\% \\
CDPP & 中证500/5日 & -37.82\% & -29.46\% & -34.77\% & -29.25\% & -31.75\% & -32.10\% \\
CDPDP-4 & 中证500/5日 & -37.85\% & -41.13\% & -35.97\% & -33.30\% & -26.29\% & -32.10\% \\
LowAmpMomentum & 中证500/5日 & -34.57\% & -34.06\% & -35.06\% & -33.62\% & -26.92\% & -32.10\% \\
Skewness & 中证500/5日 & -36.02\% & -38.14\% & -36.62\% & -32.67\% & -30.09\% & -32.10\% \\
XGBoost & 中证500/5日 & -32.17\% & -32.70\% & -32.71\% & -32.84\% & -32.13\% & -32.10\% \\
\bottomrule
\end{tabular}\end{table}
这一参照的100个多头组中，\ReferenceNegativeGroups 个2018累计收益为负。
多组同时亏损时，应先看同池等权的下跌幅度，再判断G10是否改善了持股结果。

\clearpage
\subsection{训练期选定组合的十分组收益}
每行十组都来自同一个股票池和同一个期限，没有拼接不同池的组别。

\begin{table}[H]\centering\normalsize
\caption{训练期所选组合：2018年G1至G5多头收益}
\setlength{\tabcolsep}{3pt}
\begin{tabular}{llrrrrr}\toprule
策略 & 股票池/期限 & G1 & G2 & G3 & G4 & G5\\\midrule
SmartMoney & 沪深300/5日 & -41.92\% & -44.48\% & -39.44\% & -41.68\% & -30.32\% \\
CPV & 沪深300/20日 & -35.01\% & -31.06\% & -29.36\% & -33.97\% & -29.51\% \\
CDPDP$^{\prime}$ & 沪深300/5日 & -38.03\% & -33.01\% & -38.90\% & -36.63\% & -32.05\% \\
IdealAmplitude & 沪深300/20日 & -37.61\% & -30.27\% & -31.98\% & -31.42\% & -30.28\% \\
IdealReversal & 中证500/10日 & -39.63\% & -40.57\% & -39.77\% & -40.35\% & -29.09\% \\
CDPP & 沪深300/20日 & -43.97\% & -36.70\% & -27.11\% & -32.63\% & -32.41\% \\
CDPDP-4 & 沪深300/20日 & -34.95\% & -31.87\% & -26.51\% & -29.72\% & -30.97\% \\
LowAmpMomentum & 沪深300/5日 & -33.46\% & -27.57\% & -36.17\% & -33.57\% & -29.98\% \\
Skewness & 沪深300/20日 & -38.67\% & -38.53\% & -31.14\% & -31.01\% & -31.08\% \\
XGBoost & 沪深300/20日 & -41.15\% & -40.13\% & -32.23\% & -27.66\% & -26.02\% \\
\bottomrule
\end{tabular}\end{table}
\begin{table}[H]\centering\normalsize
\caption{训练期所选组合：2018年G6至G10与全体等权收益}
\setlength{\tabcolsep}{3pt}
\begin{tabular}{llrrrrrr}\toprule
策略 & 股票池/期限 & G6 & G7 & G8 & G9 & G10 & 池等权\\\midrule
SmartMoney & 沪深300/5日 & -29.40\% & -22.54\% & -26.99\% & -34.58\% & -34.79\% & -27.15\% \\
CPV & 沪深300/20日 & -28.26\% & -26.55\% & -30.04\% & -22.72\% & -22.94\% & -27.01\% \\
CDPDP$^{\prime}$ & 沪深300/5日 & -36.85\% & -27.31\% & -28.39\% & -28.34\% & -24.95\% & -27.15\% \\
IdealAmplitude & 沪深300/20日 & -26.44\% & -22.04\% & -28.85\% & -26.40\% & -24.34\% & -27.01\% \\
IdealReversal & 中证500/10日 & -31.46\% & -33.09\% & -34.89\% & -35.03\% & -31.29\% & -32.09\% \\
CDPP & 沪深300/20日 & -29.41\% & -27.64\% & -17.47\% & -18.88\% & -19.44\% & -27.01\% \\
CDPDP-4 & 沪深300/20日 & -31.42\% & -27.83\% & -22.21\% & -25.11\% & -25.47\% & -27.01\% \\
LowAmpMomentum & 沪深300/5日 & -30.33\% & -25.95\% & -30.08\% & -29.98\% & -25.35\% & -27.15\% \\
Skewness & 沪深300/20日 & -21.06\% & -28.92\% & -24.03\% & -19.21\% & -23.68\% & -27.01\% \\
XGBoost & 沪深300/20日 & -25.65\% & -23.73\% & -18.51\% & -31.28\% & -17.03\% & -27.01\% \\
\bottomrule
\end{tabular}\end{table}
2018的全部900条分组绩效以及训练期验证分组均保存在\path{results/decile_returns.csv}，
含累计、年化与夏普；图中所选十项组合及选型依据见\path{results/selected_cases.csv}。

\clearpage
\section{为什么多头亏损：市场方向与费用}
模型标签是相对股票截面平均值的收益，学习的是池内排序；正预测值不等于股票绝对收益必定为正。
G10的目标多头仓位始终为100\%，本实验没有大盘择时或动态降低净敞口规则。
因此，选出跌幅较小的股票可以获得正相对收益，但多头仍会受到下跌市场拖累。
\begin{table}[H]\centering\normalsize
\caption{2018年股票池全体等权：固定20日调仓}
\begin{tabular}{lrrr}\toprule
股票池 & 无费累计收益 & 费后累计收益 & 费后年化\\\midrule
沪深300 & -26.69\% & -27.01\% & -27.67\% \\
中证500 & -31.80\% & -32.10\% & -32.85\% \\
全部A股 & -30.21\% & -30.52\% & -31.25\% \\
\bottomrule
\end{tabular}\end{table}

\begin{table}[H]\centering\normalsize
\caption{2018年股票池等权费后月收益：固定20日调仓}
\begin{tabular}{lrrr}\toprule
月份 & 沪深300 & 中证500 & 全部A股\\\midrule
2018-01 & 3.14\% & -1.38\% & -3.09\% \\
2018-02 & -4.55\% & -3.16\% & -4.83\% \\
2018-03 & -1.18\% & 1.88\% & 5.11\% \\
2018-04 & -3.58\% & -4.17\% & -3.80\% \\
2018-05 & 0.19\% & -1.36\% & -0.55\% \\
2018-06 & -8.72\% & -9.56\% & -10.23\% \\
2018-07 & 0.56\% & 0.37\% & -0.08\% \\
2018-08 & -6.04\% & -7.41\% & -8.16\% \\
2018-09 & 1.81\% & -0.11\% & -0.61\% \\
2018-10 & -8.28\% & -10.38\% & -9.31\% \\
2018-11 & 2.03\% & 3.75\% & 5.40\% \\
2018-12 & -5.47\% & -5.41\% & -4.29\% \\
\bottomrule
\end{tabular}\end{table}
全年累计下跌不等于每天或每月都跌，月度表可直接看出反弹区间。
月收益由该月实际日收益复利得到，全年收益是月收益的连乘，不是加总。
等权基准与实际指数权重不同，不能将它的数值直接称为官方沪深300或中证500指数涨跌幅。

\clearpage
\subsection{相对市场表现与费用反事实}
中证5005日固定XGBoost在2018的多头累计收益为\ReferenceLong，同池等权为\ReferenceEW，
相对净值收益为\ReferenceRelative。市场下行会拖累多头绝对收益，但是否跑赢等权仍需逐项判断。
\ReferenceAssessment
将相同信号和成交限制下的费率设为0，多头累计收益为\ReferenceGross，
与费后相差\ReferenceFeeGap 个百分点。它是费用反事实差，费用还会改变后续资金、成交和复利路径，
不能将现金费用简单相加当作精确收益归因。

\begin{table}[H]\centering\normalsize
\caption{训练期选定组合：2018费后多头、无费多头与等权}
\setlength{\tabcolsep}{4pt}
\begin{tabular}{llrrrr}\toprule
策略 & 股票池/期限 & G10费后 & G10无费 & 池等权 & 相对收益\\\midrule
SmartMoney & 沪深300/5日 & -34.79\% & -27.42\% & -27.15\% & -10.48\% \\
CPV & 沪深300/20日 & -22.94\% & -20.48\% & -27.01\% & 5.59\% \\
CDPDP$^{\prime}$ & 沪深300/5日 & -24.95\% & -21.25\% & -27.15\% & 3.02\% \\
IdealAmplitude & 沪深300/20日 & -24.34\% & -21.68\% & -27.01\% & 3.67\% \\
IdealReversal & 中证500/10日 & -31.29\% & -27.48\% & -32.09\% & 1.18\% \\
CDPP & 沪深300/20日 & -19.44\% & -16.70\% & -27.01\% & 10.37\% \\
CDPDP-4 & 沪深300/20日 & -25.47\% & -24.27\% & -27.01\% & 2.12\% \\
LowAmpMomentum & 沪深300/5日 & -25.35\% & -23.18\% & -27.15\% & 2.46\% \\
Skewness & 沪深300/20日 & -23.68\% & -21.07\% & -27.01\% & 4.56\% \\
XGBoost & 沪深300/20日 & -17.03\% & -15.11\% & -27.01\% & 13.68\% \\
\bottomrule
\end{tabular}\end{table}

十项训练期所选策略中，\SelectedNegative 项2018费后多头累计收益为负，
\SelectedPositiveRelative 项相对净值收益大于0。应分别看市场方向、排序效果和费用拖累，
不能只由多空为正推断多头必然盈利。无费反事实仍执行真实限价和停复牌约束。

\clearpage
\resultfigure{reference_model}{固定模型参照；股票池：中证500；调仓：5交易日；2018测试}
\clearpage
\resultfigure{selected_model}{按训练期验证所选XGBoost；股票池：\BestModelPool；调仓：\BestModelPeriod；2018测试}

\clearpage
\section{手续费、结果核验与复现}
\textbf{手续费保持双边千三：每次买、卖各收成交额的0.0015，完整往返0.003。}
两条腿分别按实际成交扣费，建仓、调仓、漂移再平衡及退出均计入。
程序核验\Series 条组合账本、\Trades 笔成交的现金、持仓、净值和费用；年化按250个交易日几何复利，
夏普按实际算术日收益计算，只展示当前口径。

全部最终模型的训练标签退出日早于2018；三折训练标签均在相应验证开始前实现。
2018模型文件哈希经检查保持冻结，选图清单可仅由训练期验证结果重建。
当前\Tests 项自动测试通过，其中包含改变2018特征不影响训练结果和最终模型权重的检查。

当前唯一代码为\path{code2_new.py}，根目录报告为\path{report.pdf}，
报告源在\path{report/}，三池当前实验分别放在\path{results/CSI300/}、\path{results/CSI500/}、\path{results/AllA/}。
训练期验证与2018测试结果分开保存，汇总在\path{results/}根目录；历史代码、报告只放在本地Git中。
\noindent\attachfile[icon=Paperclip]{../code2_new.py}\quad 完整源码、内置测试及报告生成入口。\par
\texttt{python code2\_new.py --compare-universes}\par
\texttt{python code2\_new.py --report}\par
\texttt{python code2\_new.py --self-test}

\subsection*{参考文献}
\begin{thebibliography}{9}\footnotesize
\bibitem{cpv} 东吴证券．高频价量相关性，意想不到的选股因子．2020-02-23。
\bibitem{cdpdp} 东吴证券．CPV因子错位版，价量自相关性中蕴藏的选股信息．2021-03-01。
\bibitem{ky} 开源证券．独家量价因子的高频测试．2021-09-24，第7--8页。
\bibitem{amp} 开源证券．振幅因子的隐藏结构．2020-05-16，第5--7页。
\bibitem{xgb} XGBoost官方文档．Python API与参数说明．\url{https://xgboost.readthedocs.io/}；环境版本3.2.0。
\end{thebibliography}
\appendix
\clearpage
\section{九因子自然月末十组净值}
\resultfigure{HF_SmartMoney_MonthEnd}{SmartMoney自然月末回测}
\vfill
\resultfigure{HF_CPV_MonthEnd}{CPV自然月末回测}
\clearpage
\resultfigure{HF_CDPDP_MonthEnd}{CDPDP$^{\prime}$自然月末回测}
\vfill
\resultfigure{HF_IdealAmplitude_MonthEnd}{IdealAmplitude自然月末回测}
\clearpage
\resultfigure{HF_IdealReversal_MonthEnd}{IdealReversal自然月末回测}
\vfill
\resultfigure{HF_CDPP_MonthEnd}{CDPP自然月末回测}
\clearpage
\resultfigure{HF_CDPDP4_MonthEnd}{CDPDP-4自然月末回测}
\vfill
\resultfigure{HF_LowAmpMomentum_MonthEnd}{LowAmpMomentum自然月末回测}
\clearpage
\resultfigure{HF_Skewness_MonthEnd}{Skewness自然月末回测}

\end{document}
